\chapter{KESIMPULAN}
\label{chap:kesimpulan}

\section{Jawaban Pertanyaan Kesimpulan}
Pertanyaan kesimpulan pada spesifikasi tugas didasarkan pada seluruh eksperimen. Seluruh \textit{Task} 1 sampai 6 dan Bagian E sudah dikerjakan sehingga ketujuh pertanyaan dapat dijawab.

\begin{enumerate}
\item \textbf{Konferensi dengan popularitas publikasi tertinggi.} NeurIPS, dengan 25.455 makalah (1,27\% dari 1.997.244 publikasi pada 2015--2025) dan 57.800 penulis unik. NeurIPS tetap berada di peringkat pertama ketika \textit{workshop} dan acara satelit dikeluarkan dari perhitungan. Peringkat berikutnya, yaitu CVPR dan ICASSP, bergantung pada definisi: CVPR turun ke peringkat enam jika \textit{CVPR Workshops} dikeluarkan.
\item \textbf{Asosiasi konferensi paling populer dengan konferensi lain.} NeurIPS memiliki asosiasi paling banyak, tetapi bukan yang paling kuat. NeurIPS termasuk dalam 42 pasangan \textit{frequent}, 35 dari 40 \textit{itemset} berisi lima atau enam konferensi, dan menjadi tujuan 240 \textit{rules}, terbanyak dari semua konferensi. Namun, median \textit{lift} \textit{rules} yang menuju NeurIPS hanya 23,61, kedua terendah. Semakin populer sebuah konferensi, semakin rendah \textit{lift} \textit{rules} yang menuju kepadanya (Spearman $-0{,}933$), karena \textit{lift} dibagi dengan \textit{support} konferensi tujuan. Asosiasi satu-ke-satu terkuat justru dimiliki pasangan yang lebih kecil, yaitu ECOC $\Rightarrow$ OFC (\textit{lift} 80,12) dan ASRU $\Rightarrow$ INTERSPEECH (74,17) (Subbab~\ref{sec:hasil-t5-lift}).
\item \textbf{Kombinasi konferensi yang paling sering muncul bersama.} ICML dan NeurIPS (18.503 penulis, \textit{support} 1,02\%), satu-satunya pasangan yang masih \textit{frequent} pada \textit{min\_sup} 1\%, disusul ICLR dan NeurIPS (17.852) serta AAAI dan NeurIPS (15.973). \textit{Triple} teratas adalah ICLR, ICML, dan NeurIPS (10.413 penulis), dan \textit{itemset} terpanjang berisi enam konferensi, yaitu AAAI, CVPR, ECCV, ICCV, ICLR, dan NeurIPS (2.005 penulis) (Subbab~\ref{sec:hasil-t4-fp-kombinasi}).
\item \textbf{\textit{Rule} dengan \textit{lift} tertinggi.} ECCV + ACM Multimedia + NeurIPS $\Rightarrow$ AAAI + CVPR + ICCV (\textit{lift} 172,27). Dari 2.983 penulis yang menerbitkan di ECCV, ACM Multimedia, dan NeurIPS, 61,2\% juga menerbitkan di AAAI, CVPR, dan ICCV, atau 172 kali lebih sering dibanding penulis pada umumnya. \textit{Rule} ini hanya didukung 1.827 penulis, sedikit di atas ambang 1.822, dan tidak lagi lolos jika profil nama umum yang menggabungkan banyak orang dikeluarkan sehingga perlu dibaca hati-hati. \textit{Lift} tingginya terutama berasal dari \textit{consequent} tiga konferensi yang jarang muncul bersama. Di antara \textit{rule} satu-ke-satu, \textit{lift} tertinggi dimiliki ECOC $\Rightarrow$ OFC (80,12): 66,5\% penulis ECOC juga menerbitkan di OFC (Subbab~\ref{sec:hasil-t5-top10} dan \ref{sec:hasil-t5-menarik}).
\item \textbf{Perbedaan pola pada \textit{author-based} dan \textit{author-year-based}.} Ya. Kedua representasi menjawab pertanyaan yang berbeda. \textit{Author-based} menangkap riwayat publikasi selama 2015--2025, sedangkan \textit{author-year-based} hanya menangkap konferensi pada tahun yang sama. Akibatnya, \textit{support} pasangan pada \textit{author-year-based} median 0,29 kali \textit{author-based}, dan ICCV dan ECCV tidak pernah berpasangan (10.051 transaksi pada \textit{author-based}, nol pada \textit{author-year-based}) karena keduanya tercatat pada tahun yang berbeda. Pada \textit{min\_sup} 0,1\%, \textit{author-based} menghasilkan 1.387 \textit{rules}, sedangkan setiap periode \textit{author-year-based} hanya 0 sampai 60, sehingga Bagian E memakai \textit{min\_sup} 0,01\% sampai 0,05\%. Kelompok konferensi yang ditemukan tetap sama (NeurIPS--ICML--ICLR, ACL--EMNLP, CVPR--ICCV/ECCV, GLOBECOM--ICC), dan urutan pasangan teratas relatif terjaga (18 dari 20 sama, korelasi peringkat 0,892). Kelebihan \textit{author-year-based} adalah dimensi waktunya: \textit{database} ini memperlihatkan bahwa 98\% \textit{rules} pada \textit{min\_sup} 0,02\% baru muncul pada periode akhir (Subbab~\ref{sec:hasil-e-banding}).
\item \textbf{Algoritma dengan performa terbaik.} Pada keempat \textit{min\_sup} yang diuji, Apriori dan FP-Growth setara. Keduanya menghasilkan \textit{itemset} dan \textit{rules} yang identik, dan selisih median waktunya paling besar 6\% (5,41 dan 5,19 detik pada 0,1\%), lebih kecil daripada variasi antarpengulangan. Waktu keduanya setara karena transaksi DBLP pendek dan polanya dangkal. FP-Growth tetap dipilih sebagai algoritma utama karena tidak membangkitkan kandidat. Kandidat Apriori tumbuh kuadratik terhadap jumlah item \textit{frequent}: pada 0,1\% terdapat 164.451 pasangan kandidat, dan hanya 0,24\% yang \textit{frequent} (Subbab~\ref{sec:jawab-t6}).
\item \textbf{Pengetahuan atau \textit{insight} utama.} Popularitas dan asosiasi konferensi makin terpusat pada konferensi kecerdasan buatan. Porsi sepuluh konferensi teratas naik dari 7,0\% pada 2015 menjadi 11,7\% pada 2024, dan NeurIPS tumbuh 10,2 kali. Pola asosiasinya mengikuti: pasangan teratas bergeser dari GLOBECOM + ICC pada 2015--2017 ke ICLR + NeurIPS pada 2024--2025, NeurIPS terlibat pada 67\% \textit{rules} yang baru muncul, dan \textit{itemset} berisi empat konferensi atau lebih dalam satu tahun naik dari 1 menjadi 710 (\textit{min\_sup} 0,02\%). Hubungan terkuat tetap berada di dalam satu bidang dan mengarah dari konferensi kecil ke konferensi besar, sedangkan konferensi terpopuler justru memiliki \textit{lift} rendah. Kombinasi di bidang lain, misalnya GLOBECOM dan ICC, tidak menghilang tetapi porsinya mengecil. Terakhir, sebagian pola ditentukan oleh data, bukan oleh perilaku penulis: nama umum yang menggabungkan banyak orang, konferensi yang tidak diadakan setiap tahun, dan tahun yang tidak tercatat di DBLP (Subbab~\ref{sec:jawab-e}).
\end{enumerate}

\section{Keterbatasan dan Asumsi}
\begin{enumerate}
\item \textit{Venue} diambil dari segmen kedua \textit{key} sehingga \textit{workshop} dan acara pendamping ikut terhitung sebagai konferensi induknya (7,2\% makalah menurut \textit{booktitle}). Dampaknya besar pada CVPR, CHI, dan ICC (Tabel~\ref{tab:sensitivitas}).
\item \textit{Record} \texttt{inproceedings} yang \textit{key}-nya bukan \texttt{conf/} (3.161 makalah, 0,16\%) dibuang, termasuk beberapa prosiding \textit{workshop} kecil.
\item Data 2025 tidak lengkap karena konferensi dua tahunan genap dan konferensi tahunan yang belum terindeks sehingga tidak dipakai menghitung tren dan membuat transaksi \textit{author-year-based} pada periode terakhir bias.
\item Nama penulis umum tanpa sufiks disambiguasi dapat mewakili banyak orang (\textit{Yang Liu} tercatat pada 332 \textit{venue}), dan 7.573 kelompok nama yang hanya berbeda diakritik atau huruf besar-kecil tidak digabung. Akibatnya, penulis unik cenderung terhitung lebih rendah dan sebagian transaksi menjadi terlalu besar.
\item Aturan duplikat berdasarkan judul, tahun, dan \textit{venue} dapat membuang sebagian makalah berbeda yang berjudul generik (paling banyak 421 baris, 0,02\% data).
\item ICLR tidak tercatat pada DBLP untuk 2015 dan 2016 sehingga dikeluarkan dari \textit{heatmap} indeks.
\item Pengenalan acara satelit lewat kata kunci \textit{booktitle} adalah pendekatan. Beberapa acara resmi, misalnya \textit{CHI Extended Abstracts}, dapat dianggap bagian konferensi utama.
\item Model transaksi berupa himpunan tidak menyimpan jumlah makalah per \textit{venue}: 38,3\% makalah pada \textit{author-based} dan 17,4\% pada \textit{author-year-based} tidak menambah item baru (Subbab~\ref{sec:hasil-npaper}).
\item Pratinjau perbandingan pasangan hanya mencakup \textit{venue} dengan \textit{support} minimal 0,1\% pada \textit{author-based} (574 \textit{venue}) dan belum merupakan hasil \textit{mining}.
\item Popularitas hanya diukur dari aktivitas publikasi, bukan dari sitasi atau kualitas. Pengelompokan bidang konferensi pada narasi adalah pengetahuan umum, bukan atribut dari data DBLP. Pada Bab~\ref{sec:hasil-t4} dan \ref{sec:hasil-t5}, nama bidang dicocokkan dengan judul prosiding pada data DBLP.
\item Pada \textit{min\_sup} 0,1\%, 291 transaksi dengan 100 konferensi atau lebih, yang hampir seluruhnya nama umum tanpa nomor disambiguasi, menentukan lolos tidaknya sebagian \textit{itemset}. Tanpa transaksi tersebut, 109 dari 1.514 \textit{frequent itemsets} (7\%) dan tujuh dari sepuluh \textit{rule} dengan \textit{lift} tertinggi tidak lagi lolos ambang. Hasil tetap dilaporkan apa adanya.
\item Dengan \textit{min\_conf} 50\% dan \textit{support} setiap konferensi di bawah 3,2\%, setiap \textit{rule} otomatis memiliki \textit{lift} minimal sekitar 15,8. \textit{Lift} pada data ini dipakai untuk membandingkan \textit{rules}, bukan sebagai ukuran mutlak.
\item \textit{Rules} hanya ada untuk kombinasi yang lolos \textit{min\_sup}. Hubungan kuat antara dua konferensi kecil tidak terlihat.
\item Pada Bagian E, konferensi yang tidak tercatat setiap tahun (misalnya ICCV, ECCV, ASRU, SLT) dan tahun yang tidak tercatat di DBLP (misalnya ICCC pada 2024--2025) membuat sebagian \textit{rules} muncul atau hilang karena data. Hanya 67\% \textit{rules} 2024--2025 pada \textit{min\_sup} 0,02\% yang tetap ditemukan tanpa transaksi berisi 20 konferensi atau lebih, sehingga kriteria ketahanan 80\% tidak terpenuhi pada periode ini.
\item Perbandingan waktu pada \textit{Task} 6 memakai implementasi Python murni satu \textit{thread} dengan tiga pengulangan pada satu komputer. Waktu absolutnya dapat berbeda pada implementasi lain, dan selisih median beberapa persen tidak dianggap bermakna.
\end{enumerate}
